% !PRINTED 491
\weberpage{491}

Diese beiden Permutationen aber bilden keine Gruppe, weil die Wiederholung der letzteren
\[
\begin{pmatrix} 0, & 1, & 2 \\ 2, & 0, & 1 \end{pmatrix}
\]
nicht darunter vorkommt.

Für diese Functionen gilt nun der folgende Satz:

5. \gesperrt{Ist $P$ eine Permutationsgruppe von $m$ Ziffern, und sind $\alpha, \alpha_1 \dots \alpha_{m-1}$ beliebige von einander verschiedene Grössen, so giebt es rationale Functionen der $\alpha$, sogar mit rationalen Co\"efficienten, die sich nicht ändern, wenn auf die Indices von $\alpha$ eine Permutation der Gruppe $P$ angewandt wird, und die sich ändern, wenn eine nicht zu $P$ gehörige Permutation angewandt wird.}

Um ihn zu beweisen, nehmen wir eine Function
\[
\varrho = \Phi(\alpha, \alpha_1 \dots \alpha_{m-1}),
\]
wie wir sie schon im \S.\,143 betrachtet haben, die $\Pi(m)$ verschiedene Werthe erhält, wenn die $\alpha$ auf alle mögliche Arten vertauscht werden.

Bezeichnen wir die Werthe, die $\varrho$ durch Anwendung der Permutationen einer Gruppe $P$ erhält, mit
\weq{1}{\varrho, \varrho_1 \dots \varrho_{\mu-1},}
so werden nach dem Satze 2. diese Grössen nur unter einander vertauscht, wenn auf alle zugleich eine Permutation $\pi_1$ aus der Gruppe $P$ ausgeübt wird. Denn die Anwendung der Permutation $\pi_1$ auf die sämmtlichen Functionen (1) hat denselben Erfolg, wie die Anwendung der in $P_1$ enthaltenen Permutationen auf $\varrho$. Nach 2. ist aber $P_1$ mit $P$, von der Reihenfolge abgesehen, identisch.

Bedeutet also $t$ eine Variable, so bleibt die Function
\weq{2}{\psi = (t - \varrho)(t - \varrho_1) \dots (t - \varrho_{\mu-1})}
ungeändert, wenn irgend eine Permutation aus $P$ angewandt wird. Wenn nun durch eine nicht in $P$ enthaltene Permutation die Grössen (1) in
\weq{3}{\varrho', \varrho_1' \dots \varrho_{\mu-1}'}
übergehen, so kommt wenigstens $\varrho'$ sicher nicht unter den Grössen (1) vor, und folglich ist
\[
\psi_1 = (t - \varrho')(t - \varrho_1') \dots (t - \varrho_{\mu-1}')
\]
%CONT>>
